\documentclass[12pt,a4paper]{article}
\usepackage[utf8]{inputenc}
\usepackage[T1]{fontenc}
\usepackage[french]{babel}
\usepackage{geometry}
\usepackage{graphicx}
\usepackage{listings}
\usepackage{xcolor}
\usepackage{hyperref}
\usepackage{enumitem}
\usepackage{fancyhdr}
\usepackage{titlesec}
\usepackage{array}
\usepackage{longtable}
\usepackage{tcolorbox}
\usepackage{needspace}
\usepackage{etoolbox}

\geometry{margin=2.5cm}

% Empecher les sauts de page dans les blocs de code et les tcolorbox
\AtBeginEnvironment{lstlisting}{\par\needspace{6\baselineskip}}
\AtBeginEnvironment{tcolorbox}{\par\needspace{10\baselineskip}}

% Style du code
\lstset{
    language=[Sharp]C,
    basicstyle=\ttfamily\small,
    keywordstyle=\color{blue}\bfseries,
    stringstyle=\color{red},
    commentstyle=\color{gray}\itshape,
    numbers=left,
    numberstyle=\tiny\color{gray},
    frame=single,
    breaklines=true,
    tabsize=4,
    showstringspaces=false
}

% En-tête / pied de page
\pagestyle{fancy}
\fancyhf{}
\fancyhead[L]{Firewall Protocol}
\fancyhead[R]{Documentation Technique}
\fancyfoot[C]{\thepage}

\titleformat{\section}{\Large\bfseries\color{blue!70!black}}{\thesection}{1em}{}
\titleformat{\subsection}{\large\bfseries\color{blue!50!black}}{\thesubsection}{1em}{}

\begin{document}

% ============================================================
% PAGE DE GARDE
% ============================================================
\begin{titlepage}
\centering
\vspace*{3cm}
{\Huge\bfseries Firewall Protocol\par}
\vspace{1cm}
{\Large Documentation Technique\par}
\vspace{2cm}
{\large Serious Game de sensibilisation au phishing\par}
\vspace{1cm}
{\large Développé avec Unity 6 (6000.3.3f1)\par}
\vspace{0.5cm}
{\large Plateforme cible : Android\par}
\vspace{3cm}
{\large\today\par}
\end{titlepage}

\tableofcontents
\newpage

% ============================================================
% 1. PRÉSENTATION DU PROJET
% ============================================================
\section{Présentation du projet}

\subsection{Concept}
\textbf{Firewall Protocol} est un serious game mobile développé sous Unity, dont l'objectif est de sensibiliser les joueurs à la détection du phishing. Le joueur incarne un employé qui doit trier des emails en les approuvant ou rejetant via un système de swipe (glisser gauche/droite).

\subsection{Objectifs pédagogiques}
\begin{itemize}
    \item Apprendre à identifier les emails frauduleux (phishing)
    \item Reconnaître les indices visuels et textuels d'une tentative de phishing
    \item Développer un réflexe critique face aux emails suspects
    \item Sensibiliser un large public (tous âges) aux risques du phishing
\end{itemize}

\subsection{Technologies utilisées}
\begin{longtable}{|l|l|}
\hline
\textbf{Technologie} & \textbf{Usage} \\
\hline
Unity 6 (6000.3.3f1) & Moteur de jeu \\
\hline
C\# & Langage de programmation \\
\hline
Firebase Firestore & Base de données en ligne (classement) \\
\hline
Firebase Auth & Authentification anonyme \\
\hline
DOTween & Bibliothèque d'animations \\
\hline
TextMeshPro & Rendu de texte avancé \\
\hline
JSON & Format de stockage des emails \\
\hline
PlayerPrefs & Sauvegarde locale de la progression \\
\hline
\end{longtable}

% ============================================================
% 2. ARCHITECTURE DU PROJET
% ============================================================
\section{Architecture du projet}

\subsection{Organisation des fichiers}
Le projet suit une architecture modulaire avec séparation des responsabilités :

\begin{verbatim}
Assets/
+-- Scripts/
|   +-- Managers/          # Logique centrale
|   |   +-- GameManager.cs
|   |   +-- AudioManager.cs
|   |   +-- PlayerProgress.cs
|   |   +-- LeaderboardManager.cs
|   |   +-- EffectsManager.cs
|   |   +-- ShopSystem.cs
|   +-- Data/              # Modeles de donnees
|   |   +-- EmailData.cs
|   |   +-- EmailJSON.cs
|   |   +-- EmailLoader.cs
|   |   +-- LeaderboardEntry.cs
|   +-- UI/                # Interface utilisateur
|   |   +-- MainMenu.cs
|   |   +-- EmailCardUI.cs
|   |   +-- ShopUI.cs
|   |   +-- LeaderboardUI.cs
|   |   +-- PauseMenu.cs
|   |   +-- LoadingScreen.cs
|   |   +-- FeedbackPopup.cs
|   |   +-- StreakUI.cs
|   |   +-- PowerUpBarUI.cs
|   |   +-- EndScreenAnimator.cs
|   +-- Environment/       # Elements du decor
|       +-- ApartmentScreen.cs
|       +-- ApartmentZone.cs
+-- Audio/
|   +-- Music/             # Musiques de fond
|   +-- SFX/               # Effets sonores
+-- Sprites/               # Images et sprites
+-- Resources/
|   +-- emails.json        # Base de donnees des emails
+-- Scenes/
    +-- SampleScene.unity  # Scene principale
\end{verbatim}

\subsection{Design Patterns utilisés}

\subsubsection{Singleton}
La plupart des managers utilisent le pattern Singleton pour garantir une instance unique accessible globalement :

\begin{lstlisting}
public static GameManager Instance;

void Awake()
{
    if (Instance == null)
        Instance = this;
    else
        Destroy(gameObject);
}
\end{lstlisting}

Classes utilisant ce pattern : \texttt{GameManager}, \texttt{AudioManager}, \texttt{PlayerProgress}, \texttt{LeaderboardManager}, \texttt{EffectsManager}, \texttt{MainMenu}, \texttt{EmailLoader}.

\subsubsection{Observer (callbacks)}
La communication avec Firebase utilise des callbacks asynchrones :

\needspace{15\baselineskip}
\begin{lstlisting}
public void GetTopScores(int limit, Action<List<LeaderboardEntry>> callback)
{
    db.Collection(COLLECTION_NAME)
        .OrderByDescending("bestScore")
        .Limit(limit)
        .GetSnapshotAsync()
        .ContinueWithOnMainThread(task =>
        {
            // ... traitement ...
            callback?.Invoke(entries);
        });
}
\end{lstlisting}

% ============================================================
% 3. SCRIPTS PRINCIPAUX
% ============================================================
\section{Scripts principaux}

\subsection{GameManager.cs — Logique de jeu}

Le \texttt{GameManager} est le chef d'orchestre du jeu. Il gère :

\begin{itemize}
    \item L'initialisation des parties (chargement des emails selon le jour)
    \item Le traitement des décisions du joueur (approuver/rejeter un email)
    \item Le calcul du score, des coins et de l'intégrité
    \item La gestion des power-ups (bouclier, indice, skip)
    \item Les conditions de victoire et de game over
    \item La navigation entre les écrans (appartement, menu)
\end{itemize}

\begin{tcolorbox}[title=Flux de jeu principal]
\begin{enumerate}
    \item \texttt{InitialiserPartie()} : charge les emails du jour, initialise les stats
    \item \texttt{ChargerEmailSuivant()} : affiche l'email suivant
    \item \texttt{TraiterDecisionAvecRetour()} : évalue la réponse du joueur
    \item Si correct : +score, +coins, +streak, email suivant
    \item Si incorrect : -intégrité, feedback explicatif, vérification game over
    \item \texttt{FinDeJournee()} : calcul des bonus, sauvegarde, soumission au leaderboard
\end{enumerate}
\end{tcolorbox}

\subsubsection{Gestion du bouclier}
Quand le joueur se trompe et possède un bouclier actif, les dégâts sont réduits de 50\% et le bouclier est consommé. Le SFX joué est celui du shield break (et non le SFX d'erreur classique) pour éviter la superposition de sons :

\needspace{6\baselineskip}
\begin{lstlisting}
if (shieldUsed)
    AudioManager.Instance.PlayShieldBreak();
else
    AudioManager.Instance.PlayWrong();
\end{lstlisting}

\subsection{PlayerProgress.cs — Progression et sauvegarde}

Gère toute la progression du joueur avec sauvegarde automatique via \texttt{PlayerPrefs} au format JSON.

\begin{itemize}
    \item \textbf{Progression} : jour actuel, plus haut jour atteint, coins
    \item \textbf{Statistiques} : emails traités, bonnes/mauvaises réponses, parties jouées
    \item \textbf{Boutique} : items possédés, vies, indices, skips
    \item \textbf{Streak} : série de bonnes réponses consécutives
\end{itemize}

\subsubsection{Système de niveaux (DayConfig)}
Chaque jour a une configuration dynamique calculée par \texttt{GetDayConfig(int day)} :

\begin{longtable}{|l|l|l|l|}
\hline
\textbf{Jours} & \textbf{Emails} & \textbf{Intégrité} & \textbf{Répartition} \\
\hline
1--2 & 6--7 & 100\% & 100\% facile \\
\hline
3--4 & 8--9 & 90--95\% & 60\% facile, 40\% moyen \\
\hline
5--6 & 10--11 & 80--85\% & 30\% facile, 50\% moyen, 20\% difficile \\
\hline
7--8 & 12--13 & 70--75\% & 10\% facile, 40\% moyen, 40\% difficile, 10\% expert \\
\hline
9+ & 14--15 & 60\% & 5\% facile, 25\% moyen, 40\% difficile, 30\% expert \\
\hline
\end{longtable}

Le multiplicateur de récompense augmente de 20\% par jour ($1\times$, $1.2\times$, $1.4\times$...).

\subsection{EmailLoader.cs — Chargement des emails}

Charge la base de données d'emails depuis un fichier JSON dans \texttt{Resources/}. Sélectionne les emails pour chaque journée selon la difficulté configurée, en utilisant l'algorithme de mélange Fisher-Yates.

\begin{tcolorbox}[title=Processus de sélection des emails]
\begin{enumerate}
    \item Récupère la configuration du jour (\texttt{DayConfig})
    \item Filtre les emails par difficulté (facile, moyen, difficile, expert)
    \item Mélange chaque groupe aléatoirement
    \item Sélectionne le nombre requis de chaque difficulté
    \item Mélange final pour éviter le regroupement par difficulté
\end{enumerate}
\end{tcolorbox}

\subsection{AudioManager.cs — Gestion audio}

Singleton gérant la musique de fond et les effets sonores avec deux \texttt{AudioSource} séparés.

\begin{itemize}
    \item \textbf{Musiques} : menu, gameplay, appartement — avec crossfade via DOTween
    \item \textbf{SFX} : bonne/mauvaise réponse, shield break, streak break, swipe, achat, victoire, game over
    \item \textbf{Volume} : contrôle via slider dans le menu pause, sauvegardé dans PlayerPrefs
    \item \textbf{Mute} : toggle avec changement d'icône (on/off)
\end{itemize}

Certains SFX ont des multiplicateurs de volume personnalisés :
\begin{itemize}
    \item Swipe email : $\times 4$ (pour être audible pendant le gameplay)
    \item Shield break : $\times 3$ (pour se démarquer)
    \item Streak break : $\times 7$ (pour l'impact)
    \item Skip : $\times 100$ avec pitch $\times 3$ (effet accéléré)
\end{itemize}

Le crossfade entre musiques utilise DOTween :
\needspace{10\baselineskip}
\begin{lstlisting}
musicSource.DOFade(0f, musicFadeDuration).OnComplete(() =>
{
    musicSource.clip = clip;
    musicSource.loop = true;
    musicSource.Play();
    musicSource.DOFade(isMuted ? 0f : musicVolume, musicFadeDuration);
});
\end{lstlisting}

\subsection{LeaderboardManager.cs — Classement en ligne}

Gère la communication avec Firebase Firestore pour le classement.

\begin{itemize}
    \item \textbf{SubmitScore()} : envoie le score si meilleur que l'existant
    \item \textbf{GetTopScores()} : récupère le top 8 trié par score décroissant
    \item \textbf{GetPlayerRank()} : calcule le rang du joueur
    \item \textbf{IsPseudoTaken()} : vérifie l'unicité du pseudo avant inscription
\end{itemize}

Structure d'un document Firestore :
\needspace{8\baselineskip}
\begin{verbatim}
leaderboard/{pseudo}
+-- pseudo: string
+-- bestScore: int
+-- highestDay: int
+-- timestamp: ServerTimestamp
\end{verbatim}

\subsection{LoadingScreen.cs — Écran de chargement}

Affiche un écran de chargement animé au lancement du jeu :
\begin{itemize}
    \item Fond noir avec la boîte aux lettres au centre
    \item Enveloppes animées qui tombent vers la boîte (spawn en boucle avec DOTween)
    \item Texte "Chargement..." avec animation des points
    \item Disparition en fade-out après un temps minimum (3 secondes)
    \item La musique du menu ne démarre qu'après la fin du chargement
\end{itemize}

% ============================================================
% 4. INTERFACE UTILISATEUR
% ============================================================
\section{Interface utilisateur}

\subsection{Navigation entre les écrans}

\needspace{10\baselineskip}
\noindent
\begin{minipage}{\textwidth}
\begin{verbatim}
Loading Screen --> Menu Principal --> Appartement --> Gameplay
                        |                 |              |
                        |                 |       Victoire/Game Over
                        |                 |              |
                        +-----------------+--------------+
                        (retour au menu ou a l'appartement)
\end{verbatim}
\end{minipage}

\subsection{MainMenu.cs — Menu principal}
\begin{itemize}
    \item Affiche le pseudo connecté, le jour actuel et les cryptos
    \item Bouton "Continuer" : reprend la partie (→ appartement)
    \item Bouton "Nouvelle partie" / "Recommencer" : repart du jour 1
    \item Bouton "Classement" : ouvre le leaderboard (demande un pseudo si pas encore défini)
    \item Vérification de l'unicité du pseudo via Firebase avant validation
    \item Animations d'entrée : fade-in du canvas, scale du logo, apparition séquentielle des boutons
\end{itemize}

\subsection{EmailCardUI.cs — Carte email}
Interface de swipe pour approuver/rejeter les emails. Le joueur glisse la carte à gauche (rejeter) ou à droite (approuver). Un seuil de distance est nécessaire pour valider l'action.

\subsection{ShopUI.cs — Boutique}
Permet d'acheter des power-ups avec les cryptos gagnés :
\begin{itemize}
    \item Bouclier : réduit les dégâts de 50\% (consommable)
    \item Indice : révèle si l'email est frauduleux ou légitime
    \item Skip : passe un email sans pénalité
    \item Vie supplémentaire : restaure 25\% d'intégrité en cas de game over
\end{itemize}

\subsection{LeaderboardUI.cs — Classement}
Affiche le top 8 des joueurs avec pseudo, score et jour maximum. Le texte "Chargement..." s'affiche pendant la requête Firebase. Le rang du joueur est affiché en haut du panneau.

\subsection{PauseMenu.cs — Menu pause}
\begin{itemize}
    \item Bouton reprendre (ferme le menu)
    \item Bouton mute (toggle son on/off avec changement d'icône)
    \item Bouton quitter
    \item Slider de volume connecté à l'AudioManager
\end{itemize}

% ============================================================
% 5. EFFETS VISUELS ET SONORES
% ============================================================
\section{Effets visuels et sonores}

\subsection{Effets visuels}
\begin{itemize}
    \item \textbf{GlitchEffect} : effet de glitch visuel sur mauvaise réponse (probabilité 40\%) et game over (version intense)
    \item \textbf{ConfettiEffect} : confettis sur victoire
    \item \textbf{ShieldBreakEffect} : animation visuelle de destruction du bouclier
    \item \textbf{FloatingTextEffect} : texte flottant "+score" ou "-dégâts"
    \item \textbf{StreakUI} : affichage de la série de bonnes réponses
    \item \textbf{EndScreenAnimator} : animation des écrans de fin (victoire/game over)
\end{itemize}

\subsection{Musiques}
Trois ambiances musicales avec transitions en crossfade :
\begin{itemize}
    \item Menu principal : ambiance calme d'accueil
    \item Gameplay : musique dynamique pendant le tri des emails
    \item Appartement : ambiance détendue entre les journées
\end{itemize}

\subsection{Effets sonores}
\begin{longtable}{|l|l|}
\hline
\textbf{Événement} & \textbf{SFX} \\
\hline
Bonne réponse & Validation positive \\
\hline
Mauvaise réponse & Buzzer d'erreur \\
\hline
Shield break & Son de bouclier brisé \\
\hline
Streak break & Son de série perdue (volume $\times 7$) \\
\hline
Swipe email & Son de glissement (volume $\times 4$) \\
\hline
Achat boutique & Son de caisse \\
\hline
Victoire & Fanfare \\
\hline
Game over & Son de défaite \\
\hline
\end{longtable}

% ============================================================
% 6. DONNÉES
% ============================================================
\section{Données}

\subsection{Structure d'un email (EmailData)}
\needspace{14\baselineskip}
\begin{lstlisting}
public class EmailData
{
    public string expediteurNom;      // Nom de l'expéditeur
    public string expediteurEmail;    // Adresse email
    public string objet;             // Objet du mail
    public string corpsDuMessage;    // Corps du message
    public bool estFrauduleux;       // Phishing ou non
    public string explicationErreur; // Explication pedagogique
    public int pointsSiCorrect;      // Points gagnes
    public int degatsIntegrite;      // Degats si erreur
}
\end{lstlisting}

\subsection{Fichier JSON (emails.json)}
Les emails sont stockés dans \texttt{Resources/emails.json} avec 4 niveaux de difficulté :
\begin{itemize}
    \item \textbf{Facile} : indices évidents (fautes grossières, adresses suspectes)
    \item \textbf{Moyen} : indices plus subtils (urgence artificielle, liens trompeurs)
    \item \textbf{Difficile} : très réaliste, peu d'indices
    \item \textbf{Expert} : quasiment indiscernable d'un vrai email
\end{itemize}

\subsection{Sauvegarde locale (PlayerPrefs)}
La progression est sérialisée en JSON et stockée via \texttt{PlayerPrefs} :
\begin{itemize}
    \item Clé \texttt{PlayerProgress} : toutes les données de progression
    \item Clé \texttt{MusicVolume} / \texttt{SfxVolume} : volumes audio
    \item Clé \texttt{IsMuted} : état du mute
    \item Clé \texttt{PlayerPseudo} : pseudo du joueur
\end{itemize}

% ============================================================
% 7. DÉPLOIEMENT
% ============================================================
\section{Déploiement}

\subsection{Build Android}
\begin{itemize}
    \item Plateforme : Android (APK)
    \item Target API Level : 33+ (Android 13)
    \item Minimum API Level : 22 (Android 5.1)
    \item Rendering : Universal Render Pipeline (URP)
    \item Canvas Scaler : Scale With Screen Size, référence 1080×1920, match 0.5
\end{itemize}

\subsection{Prérequis}
\begin{itemize}
    \item Unity 6 avec module Android Build Support
    \item Android SDK \& NDK Tools
    \item OpenJDK (installé avec Unity)
    \item Compte Google Play (25\$ one-time) pour la distribution
\end{itemize}

\subsection{Évolutions prévues}
\begin{itemize}
    \item Migration de la base d'emails sur Firebase pour mise à jour en temps réel sans republier l'APK
    \item Ajout de nouveaux emails basés sur les arnaques actuelles
    \item Version B2B avec emails personnalisables par entreprise
\end{itemize}

% ============================================================
% 8. DIAGRAMME DE CLASSES SIMPLIFIÉ
% ============================================================
\section{Diagramme de classes simplifié}

\needspace{30\baselineskip}
\begin{verbatim}
+-------------------+     +------------------+
|   GameManager     |---->|  EmailLoader     |
|                   |     |  - emails.json   |
| - integrite       |     +------------------+
| - score           |
| - coins           |     +------------------+
|                   |---->| PlayerProgress   |
| TraiterDecision() |     | - currentDay     |
| FinDeJournee()    |     | - coins          |
| GameOver()        |     | - ownedItems     |
+-------------------+     | Save() / Load()  |
        |                  +------------------+
        |
        v                  +------------------+
+-------------------+     |  AudioManager    |
|   EmailCardUI     |     | - musicSource    |
|                   |     | - sfxSource      |
| Swipe gauche/     |     | PlayMusic()      |
| droite            |     | PlaySfx()        |
+-------------------+     +------------------+

+-------------------+     +------------------+
| LeaderboardMgr    |---->|  Firebase        |
|                   |     |  Firestore       |
| SubmitScore()     |     +------------------+
| GetTopScores()    |
| IsPseudoTaken()   |     +------------------+
+-------------------+     |   ShopSystem     |
                          | - items[]        |
+-------------------+     | Purchase()       |
|  EffectsManager   |     +------------------+
| - GlitchEffect    |
| - ConfettiEffect  |
| - ShieldBreak     |
+-------------------+
\end{verbatim}

\end{document}
